\documentclass[12pt]{article}
\usepackage[italian]{babel}
\usepackage{graphicx}
\usepackage{hyperref}
\usepackage{setspace}

\title{Documentazione del Road Trip Generator}
\author{Tommaso Fachin}
\date{\today}

\begin{document}

\maketitle
\tableofcontents
\newpage

\onehalfspacing

\section{Introduzione}
Questo documento descrive l'analisi, la progettazione e l'implementazione di un sistema di tipo \textbf{agentico} in grado di generare automaticamente un itinerario di viaggio (road trip) sulla base di specifiche fornite dall'utente.



\section{Idea del Progetto}
L'idea iniziale nasce dall'esigenza di creare un assistente in grado di automatizzare la pianificazione di un viaggio su strada.  

\section{Obiettivi}
L'obiettivo del progetto è realizzare un agente intelligente capace di:
\begin{itemize}
    \item interpretare richieste in linguaggio naturale;
    \item estrarre informazioni rilevanti (città, preferenze, durata del viaggio, vincoli);
    \item interrogare API esterne per ottenere dati reali;
    \item generare un itinerario coerente, realistico e personalizzato.
\end{itemize}

\section {Specifiche dell' utente}
L'utente fornisce alcune specifiche, come:
\begin{itemize}
    \item città di partenza e arrivo
    \item durata del viaggio
    \item eventuali preferenze di viaggio (veloce, itinerante)
    \item nuemro massimo di ore di viaggio giornalire
    \item eventuali tappe obbligatorie
\end{itemize}

\section{Risposta dell'agente}
L'agente elabora queste informazioni e produce:
\begin{itemize}
    \item un percorso ottimizzato
    \item una lista di tappe consigliate
    \item punti di interesse lungo il tragitto
    \item ristoranti, panorami o attrazioni rilevanti
    \item tempi di percorrenza e distanze.
\end{itemize}

\section{Architettura del Sistema}
L'architettura del sistema è composta da diversi moduli principali:
\begin{itemize}
    \item \textbf{Agente LLM}: responsabile dell'interpretazione delle richieste e della generazione dell'output.
    \item \textbf{Modulo di "Routing"}: utilizza API di mappe per calcolare percorsi e distanze.
    \item \textbf{Modulo Punti di interesse}: interroga servizi esterni per ottenere punti di interesse.
    \item \textbf{Modulo Meteo}: opzionale, per integrare condizioni meteo lungo il percorso.
    \item \textbf{Interfaccia Utente}: futura interfaccia web (opure CLI).
\end{itemize}

\section{Agenti e Linguaggi utilizzati}

Il sistema utilizza un agente basato su modello LLM eseguito in locale tramite Ollama, utilizzando il modello \textbf{Mistral}.
\newline
Per la gestione dell'agente e delle sue capacità viene utilizzato il framework \textbf{PydanticAI}, che permette di definire in modo chiaro il comportamento dell'agente, le istruzioni operative e l'integrazione con funzioni esterne (API, moduli Python, ecc.).
\newline
Il progetto è sviluppato principalmente in \textbf{Python}.

\section{API Utilizzate}
DA SCRIVERE STRADA FACENDO


\section{Funzionalità Previste}
Il sistema dovrà essere in grado di:
\begin{itemize}
    \item generare un itinerario completo;
    \item adattarsi alle preferenze dell'utente;
    \item proporre alternative;
    \item gestire errori e richieste ambigue;
    \item fornire output sia testuale che strutturato (JSON).
\end{itemize}

\section{Conclusioni}

Il progetto rappresenta un esempio concreto di utilizzo di agenti AI per automatizzare compiti complessi.  
La documentazione verrà ampliata man mano che lo sviluppo procederà.

\end{document}
